\section{Dynamic feedback in the audit system}
\label{sec:results-dynamic}

The static scenarios hold enforcement parameters fixed across periods.
The simulator also supports two forms of dynamic feedback, specified
by equations~\eqref{eq:coeff}--\eqref{eq:rep}. The first multiplies
each firm's perceived reputational cost by $(1 + \varepsilon)$ after
every catch, with no decay between catches, so the accumulated
penalty grows exponentially in the cumulative number of past
violations. The second adds an increment $\Delta$ to the firm's
audit-targeting coefficient on a catch and decays the coefficient
back to its baseline at rate $\delta$ in subsequent periods, so the
elevated deterrence is bounded and dissipates over time. The two are
stylised representations of distinct deterrence intuitions: permanent
reputational memory in the first case, and time-limited regulatory
attention in the second. To avoid relying on extreme parameter values,
we identify the weakest setting of each mechanism that still produces
the characteristic behaviour and report the per-step trajectory at
that setting.

\subsection{Scenario 4 --- Reputation Ratchet}

A sweep over the reputation escalation factor $\varepsilon$ between
0 and 1.5 (30 seeds, 50 steps, audit channel off) shows mean
final-step compliance fixed at the supply floor of 25\% up to
$\varepsilon = 0.2$, rising sharply through the interval
$\varepsilon \in [0.3, 0.7]$, and saturating above
$\varepsilon \approx 1.0$. Compliance crosses 90\% at
$\varepsilon \approx 0.6$ and 95\% at $\varepsilon \approx 0.7$.
We adopt $\varepsilon^* = 0.7$ as the weakest setting reaching 95\%
and report the per-step trajectory at this value over 200 Monte
Carlo runs and 80 simulation steps (Figure~\ref{fig:ratchet}). Mean
compliance crosses 90\% near step 46, reaches 98\% by step 60, and
converges to \textbf{99.9\%} (SD\,=\,0.6\%) by step 80, at an audit
rate of 20.0\%. The trajectory is monotone because reputational
costs accumulate without decay: a firm caught enough times never
returns to a state in which violation is again the rational choice.

\begin{figure}[htbp]
\centering
\begin{tikzpicture}
\begin{axis}[
    width=0.85\textwidth,
    height=6cm,
    xlabel={Step},
    ylabel={Compliance rate},
    ymin=0.2, ymax=1.05,
    xmin=1, xmax=80,
    ytick={0, 0.2, 0.4, 0.6, 0.8, 1.0},
    yticklabels={0\%, 20\%, 40\%, 60\%, 80\%, 100\%},
    grid=major,
    grid style={gray!30},
    legend pos=south east,
]
\addplot[name path=upper, draw=none, forget plot]
    table[x=step, y=upper, col sep=comma]{data/compliance_ratchet.csv};
\addplot[name path=lower, draw=none, forget plot]
    table[x=step, y=lower, col sep=comma]{data/compliance_ratchet.csv};
\addplot[blue!25, fill opacity=0.6] fill between[of=upper and lower];
\addplot[blue, thick]
    table[x=step, y=mean, col sep=comma]{data/compliance_ratchet.csv};
\addlegendentry{Mean $\pm$ 1 s.d.\ (200 runs)}
\end{axis}
\end{tikzpicture}
\caption{Scenario 4 (Reputation Ratchet) per-step compliance at
  $\varepsilon^* = 0.7$, the weakest reputation escalation that
  converges to full compliance ($Q=5$, $N=20$, $\pi_0 = 0.20$,
  audit-channel feedback off). Compliance rises monotonically as
  repeated violations compound each lab's perceived reputational
  cost above the deterrence threshold, reaching 99.9\% by step 80.
  Shaded band is mean $\pm$ 1 standard deviation across 200 Monte
  Carlo runs.}
\label{fig:ratchet}
\end{figure}

\subsection{Scenario 5 --- Enforcement Cycles}

A joint sweep over audit escalation $\Delta \in \{0.5, 1.0, 1.5,
2.0, 3.0\}$ and decay rate $\delta \in \{0.05, 0.10, 0.20, 0.40,
0.60\}$ (reputation off, signal-dependent auditing on, 30 seeds, 60
steps) yields three regions in the parameter grid. At low decay
($\delta \leq 0.10$) the elevated audit coefficient persists long
enough that mean compliance settles above 85\%. At high decay
($\delta \geq 0.40$) deterrence dissipates between successive catches
and mean compliance collapses toward the supply floor. The
intermediate band around $\delta = 0.20$ produces visible cycles
while holding mean compliance well above the floor. We adopt
$\Delta^* = 1.0$, $\delta^* = 0.20$ as the weakest setting in this
band and report the per-step trajectory over 200 Monte Carlo runs
(Figure~\ref{fig:oscillation}). Mean compliance is \textbf{77.5\%}
(SD\,=\,0.9\% across seeds) with a five- to six-step cycle period at
an audit rate of 19.1\%. This is close to Strict Enforcement's
88.1\% compliance, achieved at roughly two-thirds the inspection
burden (19.1\% vs.\ 28.4\%).

\begin{figure}[htbp]
\centering
\begin{tikzpicture}
\begin{axis}[
    width=0.85\textwidth,
    height=6cm,
    xlabel={Step},
    ylabel={Compliance rate},
    ymin=0.2, ymax=1.0,
    xmin=1, xmax=60,
    ytick={0, 0.2, 0.4, 0.6, 0.8, 1.0},
    yticklabels={0\%, 20\%, 40\%, 60\%, 80\%, 100\%},
    grid=major,
    grid style={gray!30},
    legend pos=south east,
]
\addplot[name path=upper, draw=none, forget plot]
    table[x=step, y=upper, col sep=comma]{data/compliance_oscillation.csv};
\addplot[name path=lower, draw=none, forget plot]
    table[x=step, y=lower, col sep=comma]{data/compliance_oscillation.csv};
\addplot[red!25, fill opacity=0.6] fill between[of=upper and lower];
\addplot[red, thick]
    table[x=step, y=mean, col sep=comma]{data/compliance_oscillation.csv};
\addlegendentry{Mean $\pm$ 1 s.d.\ (200 runs)}
\end{axis}
\end{tikzpicture}
\caption{Scenario 5 (Enforcement Cycles) per-step compliance at the
  weakest audit-escalation setting that holds mean compliance well
  above the supply floor, $\Delta^* = 1.0$, $\delta^* = 0.20$
  ($Q=5$, $N=20$, $\pi_0 = 0.05$, signal-dependent auditing on,
  reputation off). Compliance oscillates around 77.5\% as
  audit-coefficient decay erodes deterrence between catches; the
  cycle period is roughly five to six steps. Shaded band is mean
  $\pm$ 1 standard deviation across 200 Monte Carlo runs.}
\label{fig:oscillation}
\end{figure}

\paragraph{Comparison.}

The two mechanisms produce qualitatively different equilibria because
their memory structures differ. Permanent reputational escalation
drives monotone convergence to high compliance, but rests on the
assumption that firms and observers do not discount past sanctions
over time. Audit-targeting escalation with decay is a more plausible
representation of how regulatory attention is allocated in practice;
the resulting equilibrium is an oscillating one rather than an
absorbing one. Within Scenario 5, the $(\Delta, \delta)$ surface
describes a trade-off between audit cost and mean compliance. At the
values used here, the mechanism approaches Strict Enforcement's
compliance with a noticeably lower audit burden, suggesting that the
more conservative form of dynamic feedback has policy value despite
not converging to full compliance.